
\subsection{Тестовое воздействие монохроматическим сигналом с линейно возрастающей амплитудой}

В данном пункте для оценки линейности тракта и визуализации работы фильтра в частотной области вместо гармонического колебания с фиксированной амплитудой использовался монохроматический сигнал с линейно нарастающей огибающей (Linear Amplitude, LA). 

Максимальная амплитуда огибающей была ограничена пределами динамического диапазона 9-битного АЦП (от 0 до 255 LSB), чтобы избежать переполнения разрядной сетки (клиппинга). Для того чтобы после демодуляции сигнал не оказался на нулевой частоте, частота входного воздействия была смещена на 1 МГц относительно частоты цифрового гетеродина ($f_c = f_g + 1$ МГц). Результаты моделирования представлены на рисунке \ref{fig:p5_la}.

\begin{figure}[H]
    \centering
    %\includegraphics[width=0.9\textwidth]{5.PNG}
    \caption{Временные и спектральные характеристики сигналов при воздействии сигнала с линейно возрастающей амплитудой}
    \label{fig:p5_la}
\end{figure}

\textbf{Выводы по результатам моделирования:}
Анализ временных диаграмм подтверждает корректную работу устройства: огибающая демодулированного и отфильтрованного сигнала (\texttt{Signal LA on IF = 0}) сохраняет строго линейный характер нарастания на всем временном интервале. Это свидетельствует об отсутствии амплитудных искажений и переполнений в регистрах CIC-фильтра и КИХ-корректора при работе в пределах допустимого динамического диапазона. 

Рассматривая спектр на выходе устройства (\texttt{Spectrum LA on IF = 0}), можно отчетливо наблюдать формирование частотной характеристики спроектированной системы фильтрации. Шумовая полка наглядно отрисовывает прямоугольную форму суммарной АЧХ (каскада из CIC и КИХ фильтров) с плоской вершиной и резкими спадами. Полезный тестовый сигнал в виде узкого пика корректно позиционируется ровно на частоте 1 МГц, находясь в пределах рабочей полосы пропускания, тогда как все внеполосные спектральные компоненты успешно подавлены до уровня затухания в полосе задерживания.
